\section{Mean length of a prescribed half-plane arch}\label{c:mean}
The same-side probability theorem does not control mean length.
We give a separate marked-arm proof, whose length-weighted contributions
are summable over arch heights at each prescribed gap.
It uses the calibrated finite method and its renewal consequences.

\begin{theorem}[Every-gap half-plane mean length]\label{c:arch-mean}
For two half-plane boundary ports separated by each integer $g\to\infty$,
the probability law $\kappa^{L(\gamma)}/K_g$ satisfies
\[
 L(\gamma)=g^{4/3+o(1)}\quad\hbox{in probability},\qquad
 \mathbb E L(\gamma)=g^{4/3+o(1)}.
\]
\end{theorem}
The finite calibrated estimates in \cite[Theorem 1.2 and Proposition 11.1]{Calibrated} give
$B_h\asymp h^{-1/4}$, $K_g\asymp g^{-5/4}$ with its lower bound
retained at diameter at most $Ag$ for fixed $A$, an exponential lateral
cutoff for a fixed-height strip, and total chord length mass at most
$CR^{25/12}$ in a convex box of size $R$. Their precise confined
length window is not replaced by a pointwise endpoint estimate here.
The probability inputs below are derived from this finite package
in \cite[Proposition 4.1, Lemma 4.2, Lemmas 7.2 and 7.4,
Section 7.3, and Theorem 7.1]{Renewal}.
Proposition 4.1 and Lemma 4.2 give the opening height, diameter,
and length-deficit estimates; Lemma 7.2 gives the atom bounds;
Lemma 7.4 and Section 7.3 give the censored survival test and its
shared-turn application; Theorem 7.1 gives the absolute short-path bound.

\begin{proposition}[Calibrated renewal inputs]\label{c:calibrated-inputs}
For the port-rooted irreducible law, write $Y,Z,L,D$ for integral height,
lateral displacement, length and diameter. Then $2Z$ is integral,
$\Pr(D>r)\le Cr^{-3/4}$,
$1-\mathbb E e^{-tY}\asymp t^{3/4}$, and
$1-\mathbb E e^{-tL}=t^{9/16+o(1)}$. In addition:
\begin{enumerate}
\item[(i)] For $k$ independent whole pieces,
$\sup_j\Pr(S_Y(k)=j)\le Ck^{-4/3}$ and
$\sup_w\Pr(S_{(Y,Z)}(k)=w)\le Ck^{-8/3}$.
The off-axis bridge construction gives positive mass in directional
windows; all atoms use the actual lattice.
\item[(ii)] Start two independent sequences at the same bottom port
through the same first center $v$. Retain completed pieces at heights
strictly below $s$, excluding the first piece reaching or passing $s$.
Fix left and right roles. At every pair of renewal levels
$0<d_1<d_2<s$ from the left and right histories respectively, require
that the prefixes be disjoint after $v$, leave on opposite sides of $v$,
and that every right-prefix crossing at $d_1$ lie strictly right of the
left terminal coordinate. The probability of all these requirements is
at most $C_\eta(1+s)^{-3/4+\eta}$ for each $\eta>0$.
At a shared turn this applies with $s=1$ vacuously, or with $s$
sufficiently below both adjacent arm heights that distant endpoint
adjustments do not affect the tests. For the two arms used below, extending both to the same top port
duplicates just their highest center. The resulting factor of $\kappa$
is absorbed by an absolute constant; no shape-dependent trimming is
used in this application.
\item[(iii)] For every sufficiently small fixed $\xi>0$, some $c,C>0$
give the absolute critical mass estimate
\[
 \sum_{\substack{\omega\text{ rooted at a fixed center}\\
 R\le\diam\omega\le2R,\ L(\omega)\le R^{4/3-\xi}}}
 \kappa^{L(\omega)}\le C e^{-R^c}.
\]
\end{enumerate}
\end{proposition}

\subsection{The two arms and the dyadic target}
Split an arch at a chosen highest vertex. Its two real incident edges go
downwards. Extending both arms to the directly adjacent top wall port
makes them strict downward bridges, duplicating just that vertex. Their
critical product weight therefore differs from the arch weight by a
fixed factor. We count the two shapes from a reference top port, impose
equal heights and transverse displacement difference $g$, and translate
the first bottom endpoint to its prescribed position. This entails no
additional placement sum.

The two ends of this decomposition impose different constraints. At the
top the arms share a turn, so Proposition~\ref{c:calibrated-inputs}(ii)
applies. At the bottom they start at distinct ports a distance $g$ apart;
we will prove a separate weak-order estimate there. To sum length, mark
one irreducible on one arm with its length, and reserve independent
pieces away from both ends and the mark to match height and displacement.
Adding the two choices of rewarded arm bounds the full length weight.

Put $d_*=4/3$ and $a=3/4$. For a dyadic scale $H$, let
\[
 \mathcal M_g(H)=
 \sum_{\substack{\gamma:\,0\to g\text{ a half-plane arch}\\
                  H\le h(\gamma)<2H}}
       L(\gamma)\kappa^{L(\gamma)},
\]
where $h(\gamma)$ is the integral height of the adjacent top wall just
described. The target estimate is
\begin{equation}\label{c:arch-dyadic-target}
 \mathcal M_g(H)\le C_\eta g^{1/4}H^{d_*-2a+O(\eta)}
 \qquad\bigl(H\ge g/\log^2g\bigr).
\end{equation}
The smaller height scales will be negligible by strip confinement.
Since $d_*-2a=-1/6$, summing this bound for sufficiently small $\eta$
gives total length mass $g^{1/12+o(1)}$. Division by
$K_g\gtrsim g^{-5/4}$ then gives the required mean upper bound.

\subsection{A local truncated reward}
The first ingredient bounds the cost of the piece carrying the length
weight.
\begin{lemma}[Truncated length reward]\label{c:truncated-reward}
For the strict irreducible law, as $H\to\infty$,
\[
 \mathbb E[L; Y\le 2H]\le H^{d_*-a+o(1)} .
\]
\end{lemma}
\begin{proof}
Write $\mathcal L_h$ for the total length-weighted mass of strict bridges
of height $h$ from one fixed port, with terminal port summed. Insert a
marked irreducible of height at most $2H$ between two arbitrary strict
bridges of heights at most $H$. Unique irreducible factorization gives
\[
 \left(\sum_{0\le u\le H}B_u\right)^2
       \mathbb E[L;Y\le2H]
 \le \sum_{0\le h\le4H}\mathcal L_h.
\]
Indeed the same output can have several eligible marked pieces, but
their lengths sum to at most its full length. The bridge sum on the left
is of order $H^a$.

For height $h$, restrict first to diameter at most $h\log^2h$. Translate
the source to order $h$ positions and apply the calibrated convex-box
length bound in a box of size $O(h\log^2h)$. This gives
$\mathcal L_h\le h^{13/12+o(1)}$ for the restricted bridges. The larger
diameter tail is negligible by the exponential strip cutoff and the
area bound on length. Bounded strip heights also have exponential
diameter tails: extend them a fixed amount to a sufficiently large
fixed height. Consequently
\[
 \sum_{h\le4H}\mathcal L_h\le H^{a+d_*+o(1)}.
\]
The insertion inequality now proves the claim.
\end{proof}

\subsection{Survival of the order at separated bottom ports}
The shared-turn estimate does not apply at two different starting ports.
We therefore prove the separate weak-order estimate required at the
bottom ends of the arch. This estimate uses the calibrated finite inputs
throughout, independently of the off-axis construction in
Section~\ref{c:same-side}.

\begin{lemma}[Weak order from separated ports]\label{c:bottom-survival}
Start two independent upward irreducible sequences at bottom ports a
distance $g\ge1$ apart, assigning left and right roles by their starts.
Let $\Pi_g(S)$ be the probability that their horizontal endpoint order
is preserved weakly at every common renewal height strictly below $S$.
For $S\ge2$,
\[
 \Pi_g(S)\le C(g/S)^{1/4}\sqrt{\log(2+S)}.
\]
\end{lemma}
\begin{proof}
Their strictly positive common renewal levels themselves form renewals;
write $(T,W)$ for one step, recording height and the signed change of
separation. Pairs reaching the same height split there freely. Their
renewal mass at $j$ is $V_j=B_j^2\asymp(1+j)^{-1/2}$, so the first common
return has total mass one and the pairs restart independently.
Write $b=1/2$. The transform gives
$1-\mathbb E e^{-T/r}\asymp r^{-b}$. Exchanging the two sequences preserves
the joint law of $(T,W)$ and replaces $W$ by $-W$; hence height damping
also preserves displacement symmetry. We first establish
\[
 \Pr(|W|>r)\le C r^{-b},\qquad
 \mathbb E\min(1,|W|/r)\gtrsim r^{-b}.
\]
Pairs in a height bin around \(c_0 r\), with \(c_0\) small, carry summed mass \(\gtrsim r^b\) within individual diameter \(r/3\) by tightness, hence with all \(W\) increments bounded by \(r\). Bounding this by the restricted renewal sum proves the first claim. For the second take integral height \(h\asymp r\). We need strict bridges of height $h$ with endpoint displacement of a fixed nonzero sign and size at least $ch$, each of total mass $\gtrsim h^{-1/4}$. The half-plane continuation in \cite[Section~11, Recoverable half-plane growth]{Calibrated} uses two pure inward normals with dot product $1/2$. For each fixed sufficiently small transverse tolerance, it continues from initial scale $r$ to the exact wall of the second normal at distance $S\ge Kr$, with its terminal port free. Its mass is at least $cS^{-1/4}(S/r)^{-C}$ and diameter at most $C'S$. All centers stay after the initial wall and before the terminal wall, the continuation has the prescribed relative transverse tolerance, and its joining cuts are recoverable. The constants and lower threshold for $r$ may depend on that tolerance.

Apply it first with normals $y,u$ at angle $\pi/3$, initial scale $c_*h$,
and each exact $u$-wall whose level lies in
$[(2-2\delta_*)h,(2-\delta_*)h]$, always summing its terminal ports.
For the later continuation with normals $u,y$, fix in advance a small
transverse tolerance and a sufficiently small initial-scale fraction
compatible with its target-to-initial-scale requirement. Its diameter
constant is then fixed independently of $\delta_*$. Choose $\delta_*$
small in terms of that constant, and then choose the first continuation's
tolerance and $c_*$ sufficiently smaller, also satisfying its required
scale ratio.

The first continuation gains $y$-height with slope $1/2$, so it stays in
$0<y<h$, ends with remaining height between $\delta_*h/4$ and
$2\delta_*h$, and has lateral displacement at least $ch$. Each target
wall has mass $\gtrsim h^{-1/4}$ because all scale ratios are fixed.
Extract the prefix at the lowest qualifying singly crossed $u$-cut in
this bin, using these confinement and terminal-position tests. The mass
summed over target levels is $\gtrsim h^{3/4}$, whereas all suffixes after
a given prefix cost at most $\sum_{j\le Ch}B_j\le Ch^{3/4}$. Distinct
prefixes therefore have mass $\gtrsim1$.

From each prefix apply the $u,y$ continuation with the preselected
tolerance and initial-scale fraction to reach the wall $y=h$ exactly,
again with terminal port free. It stays beyond the joining $u$-cut,
has diameter $O(\delta_*h)$ and mass $\gtrsim h^{-1/4}$. By the choice
of $\delta_*$, it preserves the nonzero lateral displacement and avoids
$y=0$. The joining cut is recovered by its first-qualifying rule. This
proves the asserted strict-bridge mass within the calibrated method.

Reflect the construction for the second bridge. The critical product mass with total transverse difference $\ge ch$ in absolute value is therefore $\gtrsim h^{-1/2}=h^{b-1}$. Bound using subadditivity of \(\min(1,|\cdot|/r)\) over common steps and renewal insertions, costing at most \(\max_{s\le h}\sum_{i+j=s}V_i V_j\) (bounded) times the asserted expectation. This proves the claim. As in the local bounds before (by tail truncation and symmetrization) these imply \(\Pr(c r\le |W|\le r)\gtrsim r^{-b}\) and \(1-|\phi_W(u)|^2\gtrsim |u|^b\) near zero for the characteristic function. Thus independent sums of \(k\) such increments have atoms \(\le C k^{-1/b}\), by lattice Fourier inversion (nonconstant support on a discrete grid).

\paragraph{Counting near-minimum times.}
The distinction between weak and strict survival for lattice increments
is essential; compare \cite[Section 2.2]{Vysotsky2014}. We use the
following direct factorization, which counts all near-minimum times.
Kill at an independent exponential height $E$ of mean $S$. The transverse
walk has symmetric subprobability step kernel
\[
 Q_S(w)=\mathbb E[e^{-T/S};W=w],\qquad
 q_0=1-\sum_wQ_S(w)\asymp S^{-b}.
\]
Let $p$ be its probability of having all partial sums at least $-g$
before death. Write $N_g$ for the number of indices, including time zero,
lying within
$g$ of its global minimum. Then $\mathbb E N_g=p^2/q_0$. Here is a count that
also handles ties. If $a_j(g)$ is the subprobability mass of $j$ steps
whose every partial sum is at least $-g$, then
$p=q_0\sum_{j\ge0}a_j(g)$. At a candidate time with $i$ steps before
it and $j$ after it, reverse and negate the past increments. Symmetry
and independence give mass $q_0a_i(g)a_j(g)$ for the candidate to lie
within $g$ of the minimum. Summing in $i,j$ gives $p^2/q_0$.
Every qualifying candidate is counted; no unique minimum is selected.

\paragraph{Bounding the occupation count.}
The killed kernel is dominated by the original displacement law, including
after convolution:
\[
 Q_S^{*k}(w)\le\Pr(W_1+\cdots+W_k=w).
\]
Thus the atom estimate above bounds the expected number $L_A$ of visits
to any interval $A$ of width $2g$, uniformly in the starting position, by
\[
 \mathbb E L_A\le 1+\sum_{k\ge1}\min\{1,Cgk^{-1/b}\}
 \le Cg^b.
\]
Markov's inequality and restarting after successive visit-count thresholds
give an exponential tail for $L_A$ at scale $g^b$. Anchor an interval of
radius $g$ at each alive index in turn and count its future visits. The
earliest near-minimum index captures all $N_g$ near-minimum visits in its
interval. We do not condition at that retrospectively chosen index:
instead union bound over all indices, whose survival probabilities sum to
$1/q_0$. Consequently
\[
 \Pr(N_g\ge t)\le
 \min\{1,(C/q_0)e^{-ct/g^b}\},\qquad
 \mathbb E N_g\le Cg^b\log(2+1/q_0).
\]
Together with the factorization this bounds $p$ by
$C(g/S)^{b/2}\sqrt{\log(2+S)}$.

Finally, preserving weak order below $S$ also preserves it below $E$ on
$\{E<S\}$. Independence of $E$ from the original sequences gives
\[
 (1-e^{-1})\Pi_g(S)\le p.
\]
Since $b/2=1/4$, this proves the lemma.
\end{proof}

\subsection{Summing the marked-arm decomposition}
\begin{proof}[Proof of Theorem~\ref{c:arch-mean}]
We now have the two different survival tests needed for an arch: the
shared-turn test at its top and Lemma~\ref{c:bottom-survival} at its bottom.
Write $b=1/2$, as in the proof of that lemma.
Use the two-arm decomposition above, whose actual arms are disjoint after
the turn. We will prove \eqref{c:arch-dyadic-target} by allocating independent
groups of pieces to these tests, the reward, and the two-coordinate match.

Use dyads \(H\le h<2H\). Those \(H<g/\log^2 g\) contribute negligibly by exponential strip aspect since one bridge has diameter of order at least \(g\). On remaining dyads diameters above \(C H\log^2 H\) cost \(O_A(H^{-A})\) for arbitrary fixed \(A\), still by strip aspect (sum bridge products at each height, also with lengths using area). Thus in discards with superpolynomial accuracy we may pay a polynomial bound for lengths.

Let \(k_i\) be the two irreducible counts, in dyads at scales \(K_i\), \(K=\max K_i\). Treat their bridge masses at specified counts as product iid laws. Discard \(k_i>H^{a+\eta}\) by height damping, with \(\eta>0\) arbitrarily small fixed (\(k_i\le2H\)). Expand the length weight by marking one piece on its arm. On the other arm also mark a large height jump \(>H^{1-\eta/(2a)}\) (upper bound by counting choices), if \(k_i\le H^{a-\eta}\); without one the height requirement costs superpolynomially little, by exponential weighting at the inverse cutoff on the truncated jumps, using \(\Pr(Y>t)\le C t^{-a}\). For larger count on that arm just designate a middle index instead.

Sort the index distances (counting from 1) of each designated piece from the top and bottom ends into dyadic scales \(I_i,J_i\) respectively. Thus \(\max(I_i,J_i)\asymp K_i\). The sum of index choices and the associated length or large-jump costs have budget, on the two sides combined,
\[
 C_\eta H^{d_*-2a+O(\eta)}
       \prod_{i=1}^2\min(I_i,J_i).
\]
Indeed for a given length piece use Lemma~\ref{c:truncated-reward}; the jump costs \(C H^{-a+O(\eta)}\), and the number of choices on each side is at most a constant times its smaller distance scale. In the middle-index case \(\min(I_i,J_i)H^{-a+\eta}\gtrsim1\) already.

Put $I=\min I_i$ and $J=\min J_i$. The following deterministic allocation makes all uses of independence explicit. Let $r_i$ be the designated index on arm $i$, whose total count is $k_i$, and write
\[
 u_i=r_i-1,\quad v_i=k_i-r_i,\quad
 u=\min_i u_i,\quad v=\min_i v_i.
\]
Take the first $\lfloor u/16\rfloor$ pieces of each arm for the top test and the last $\lfloor v/16\rfloor$ for the bottom test. Neither block contains a designated piece. They are disjoint since $u+v\le k_i-1$ for both arms. On an arm with maximal count $k_{\max}$, at least $15k_{\max}/16-1$ indices remain after these blocks and its designated piece are removed. Reserve the first $\lfloor k_{\max}/4\rfloor$ of these unused indices in their original order for the planar atom; condition on all remaining unreserved draws. The reserved draws need not be consecutive: their displacements still add and they remain independent with the original law.

If $u<16$, the top distance scale $I$ is bounded, so we omit that test and bound its probability by one; this changes the displayed cost below only by a constant. The same holds for $v<16$ and $J$. Otherwise the two testing blocks have orders $I$ and $J$. Bounded $k_{\max}$ is absorbed into constants in the atom bound. This allocation therefore covers marks next to endpoints, unequal arm lengths, and both possible designations. The factors to be multiplied are summarized here; $p'=b/(2a)=1/3$, $m_i=\min(I_i,J_i)$, and the displayed powers of $H$ allow fixed multiples of $\eta$.
\begin{center}
\begin{tabular}{@{}lll@{}}
\toprule
Piece group & Role & Upper cost \\
\midrule
Top blocks on both arms & Shared-turn order & $I^{-1}H^{O(\eta)}$ \\
Designated pieces and indices & Reward and other-arm mark & $H^{d_*-2a+O(\eta)}m_1m_2$ \\
Reserved larger-arm indices & Two-coordinate match & $K^{-2/a}$ \\
Bottom blocks on both arms & Separated-port weak order & $g^{b/2}J^{-p'}H^{O(\eta)}$ \\
\bottomrule
\end{tabular}
\end{center}
The second row includes the designated-index counting allowance. For a large-jump mark this is the actual sum of index choices. In the deterministic middle-index case its cost one is bounded by $m_iH^{-a+\eta}\gtrsim1$ on that arm, so the same combined budget applies; the factor $H^{-a+\eta}$ is not a probability bound on that unconditioned piece. Choices of total counts contribute at most $CK_1K_2$ separately.

With \(S_I=I^{1/a}H^{-C_0\eta}\) for fixed sufficiently large \(C_0\), this many pieces reach height \(S_I\) except with superpolynomial error by height damping if \(S_I\ge2\); similarly at \(S_J=J^{1/a}H^{-C_0\eta}\). Both thresholds are far below \(h\). On valid arms we thus impose Proposition~\ref{c:calibrated-inputs}(ii) up to \(S_I\) (rounded down), by the shared-turn crosscut argument in \cite[Section 7.3]{Renewal}, with the common top-port convention stated in Proposition~\ref{c:calibrated-inputs}(ii), at probability cost \(\le C_\eta I^{-1}H^{O(\eta)}\), allowing roles. At the other end impose weak order at the common cuts below \(S_J\) at cost \(\le C g^{b/2}J^{-b/(2a)}H^{O(\eta)}\), since the two disjoint prefixes to a common cut are strip crosscuts. For thresholds \(<2\) the corresponding upper costs with enlarged \(O(\eta)\) need no test. Reversed shapes are again iid by bridge symmetry, tested as if placed at the prescribed two bottom ports (actual relative placement upon imposing the increment equalities); these tests therefore do not consume information from the middle increments.



In addition the required exact height equality and transverse difference together cost \(\le C K^{-2/a}\) by the local estimate on \((Y,Z)\), reserving order \(K\) other draws of the larger count side independently of both tests and the marks (adjust constant for bounded \(K\)). Including at most \(C K_1 K_2\) pairs of counts, all count-scale factors cancel: writing \(p'=b/(2a)=1/3\),
\[
 K_1K_2 K^{-2/a}\prod_i\min(I_i,J_i)\, I^{-1}J^{-p'}
 \le C K^{-2/a}\max(I_1,I_2)\max(J_1,J_2)J^{1-p'}
 \le C.
\]
Thus including polylogarithmic bin sums proves \eqref{c:arch-dyadic-target}.
Since $d_*-2a<0$, summation and free choice of $\eta$ give total length
mass at most $g^{1/12+o(1)}$ for both arms.

The length-weighted bound is $K_g\mathbb E L\le g^{1/12+o(1)}$.
Dividing by $K_g\ge cg^{-5/4}$ proves the mean upper bound and the
long-length probability bound by Markov's inequality. For the lower
probability bound, every arch has diameter at least a fixed multiple of
$g$. Sum Proposition~\ref{c:calibrated-inputs}(iii) over dyadic diameter
bins starting there, taking a slightly smaller exponent slack to absorb
unit and endpoint conventions. The resulting mass is superpolynomially
small compared with $K_g$. Thus the short-length probability tends to
zero; it also gives $\mathbb E L\ge g^{4/3-o(1)}$. This proves
Theorem~\ref{c:arch-mean}.
\end{proof}


The following terminal-summed consequence also appears in the proof of
\cite[Theorem~8.1]{Renewal}. We include its normalization and first-moment
argument to distinguish this confined mixture from the individual-gap law.
\begin{corollary}[Confined macroscopic-gap arches]\label{c:confined-arch-mean}
Fix a half-plane boundary source and choose a sufficiently large fixed
$C$. Give each arch with terminal gap in $[h,2h]$ on one chosen side
and diameter at most $Ch$ its critical weight, then normalize over
this entire set of terminal ports. For every integer $h\to\infty$,
\[
 L(\gamma)=h^{4/3+o(1)}\quad\hbox{in probability},\qquad
 \mathbb E L(\gamma)=h^{4/3+o(1)}.
\]
\end{corollary}
\begin{proof}
The confined pointwise arch lower bound from the calibrated finite
input above, summed over the order $h$ allowed gaps,
gives normalizing mass at least $ch^{-1/4}$ once $C$ is fixed sufficiently
large. For the first moment, translate the source to order $h$ distinct
ports on the same straight side. All translated confined arches fit in
one convex pure box of diameter $O_C(h)$. Distinct sources distinguish
the translated paths, so the calibrated convex-domain length bound
\cite[Proposition 11.1]{Calibrated}, divided by the number of sources,
gives first length mass at most $C_C h^{13/12}$. Division by the
normalizer proves $\mathbb E L\le C_C h^{4/3}$ and the upper probability
tail by Markov's inequality. Every retained arch has diameter at least
a fixed multiple of $h$. Proposition~\ref{c:calibrated-inputs}(iii),
summed over the finitely many diameter bins through $Ch$, makes its
short-length mass superpolynomially small. This proves the lower
probability tail and hence the mean lower bound. The law in this
corollary sums over terminal ports before normalization; it is distinct
from the individual-gap law of Theorem~\ref{c:arch-mean}.
\end{proof}
